\setchapterpreamble[u]{\margintoc}
\chapter{From Column Space to Linear Layers}
\labch{column-space-to-layers}

\sources{Strang, \emph{Introduction to Linear Algebra}, 6th ed.\ \cite{strang2023ila},
Chapter~10; Zhang, Lipton, Li and Smola, \emph{Dive into Deep Learning}
\cite{d2l2023}, Chapters~1--5.}

The four parts of this book were assembled separately. This closing chapter says
how they meet, and states --- deliberately conservatively --- what remains to be
studied before the research question that gives the book its title can be
approached.

\section{A network is a composition of the objects already met}
\labsec{network-as-composition}

A feed-forward network alternates two kinds of map. The affine part,
\[
	Z=XW+\vect{1}\vect{b}\tp,
\]
is the batched matrix product of \refsec{batching}; the activation part applies a
fixed scalar function elementwise. Nothing in either is new. What the
composition adds is depth, and depth is what makes the resulting function
non-linear.

\begin{remark}
Stacking affine maps with no activation between them achieves nothing. A
product of matrices is a matrix, so the composition collapses to one linear map
with the same column space. The activation is the only reason a second layer is worth
having, which is a statement about \refsec{four-ways} rather than about biology.
\end{remark}

\section{Three questions, three answers already given}
\labsec{three-questions}

\paragraph{What can the model represent?}
For a single linear layer the reachable outputs are exactly \(\Col(W\tp)\)
applied to the inputs, so the rank of the weight matrix bounds what the layer can
express --- the reachability question of \refsubsec{reachability} in applied
clothing. Compressing a trained layer by truncating its singular values, as in
\refsec{low-rank}, is a deliberate reduction of that rank.

\paragraph{How is it fitted?}
By gradient descent, whose direction is justified by \refsec{steepest}, whose
step size is constrained by the curvature in \refeq{taylor2}, whose convergence
rate is governed by the condition number in \refsec{gd-quadratic}, and whose cost
is made acceptable by the association order in \refsec{backprop}.

\paragraph{Why does the fitted model generalise?}
Because training loss is a sample average of an expectation
(\refsec{lln}), because the deviation of that average shrinks like \(1/\sqrt n\)
(\refsec{clt}), and because the loss itself is a likelihood, so that penalising
the weights is placing a prior on them (\refsec{bayesian}).

\section{Where the linear picture runs out}
\labsec{limits}

Every question above was answered by treating the layer as linear and the
activation as an afterthought. That works because the analysis is local: the
Jacobian of \refsec{matrix-chain-rule} is a genuine linear map, and the whole of
Part~I applies to it at each point. It does not describe the function globally,
and it says nothing about which composition of layers should be preferred.

Two limitations are worth naming, because they are what motivates looking
further:

\begin{description}
	\item[The non-linearity is fixed.] In a conventional layer the parameters
	are the entries of \(W\) and \(\vect{b}\); the activation is chosen in
	advance and not learned. All the flexibility sits in the linear part.
	\item[Interpretation is indirect.] The learned object is a large weight
	matrix whose singular vectors carry no meaning supplied by the problem.
	Understanding what a layer has learned is a separate research effort.
\end{description}

\begin{remark}
Kolmogorov--Arnold networks are motivated by the first of these: they
parameterise the univariate functions themselves rather than the linear mixing.
Whether that reparameterisation delivers the advantages claimed for it is the
question this repository exists to study, and it is not settled by anything in
this book. \textbf{[needs verification]} --- the technical content of that
literature has not yet been read, and no claim about it is made here.
\end{remark}

\section{What to read next}
\labsec{next}

The foundations assembled here are the standard ones, and they are sufficient to
begin. The material that has not been covered, and that the next stage of these
notes should add, is:

\begin{enumerate}
	\item tensors and automatic differentiation in practice, continuing the
	notebooks already in this repository;
	\item approximation theory for univariate functions, in particular spline
	bases, which is the mathematics a Kolmogorov--Arnold parameterisation
	rests on;
	\item the Kolmogorov--Arnold representation theorem itself, in its original
	form, before any neural interpretation of it;
	\item the primary literature on these networks, read against that
	background rather than as a substitute for it.
\end{enumerate}

\marginnote{The order matters. Reading the applied literature before the
representation theorem invites accepting an analogy as a proof --- which is the
failure mode these notes are written to avoid.}

Part~I through Part~IV will be revised as that reading proceeds. The notes are
intended to be maintained rather than finished.
